\chapter{指向性アンテナを用いたWi-Fi信号強度に基づく自己位置推定の開発}
\label{cha:field_experiments}
\graphicspath{{figures/05}}

\section{はじめに}
\section{指向性アンテナを用いたWi-Fi地図作成}
Wi-Fi地図の地図構築に必要なステップは下記の通りである．
\subsection{指向性アンテナを用いたデータ収集}
\subsection{指向性アンテナを用いた強度情報の計算}
\subsection{ガウス過程回帰を用いた信号強度によるWi-Fi地図作成}
\subsection{ガウス過程回帰を用いたAoAによる角度地図作成}

\section{指向性アンテナによるWi-Fi MCL}
\subsection{Wi-Fi信号強度とWi-Fi地図に基づく尤度計算}
\subsection{AoAによる角度推定と角度地図に基づく尤度計算}
\subsection{角度尤度の適応条件}

\section{指向性アンテナによるWi-Fi MCLの精度評価実験}
\subsection{実験目的}
\label{sec:exp_purpose}

\eqref{eq:normal_distribution}はガウス分布です．
\begin{equation}
  f(x \mid \mu, \sigma^2) = \frac{1}{\sqrt{2\pi\sigma^2}} \exp\left(-\frac{(x-\mu)^2}{2\sigma^2}\right)
  \label{eq:normal_distribution}
\end{equation}



\subsection{実験環境}
\label{sec:exp_environment}

実験環境を説明します．


\subsection{実験結果}
\label{sec:exp_results}



\section{考察}
\label{sec:exp_discussion}



考察します．
\section{おわりに}